%!TEX root = ../../exa-ma-deliverables.tex
% Sources: Exa-MA news archive, NumPEx announcements and event programmes through 18 November 2026.
\begingroup
\small
\setlength{\tabcolsep}{4pt}
\setlength{\LTpre}{6pt}
\renewcommand{\arraystretch}{1.04}
\begin{longtable}{@{}%
  >{\raggedright\arraybackslash}p{0.16\textwidth}
  >{\raggedright\arraybackslash}p{0.30\textwidth}
  >{\raggedright\arraybackslash}p{0.49\textwidth}@{}}
\caption{Selected Exa-MA and NumPEx announcements and events in 2026.}
\label{tab:exama-news-2026}\\
\toprule
\textbf{Date} & \textbf{Announcement or event} & \textbf{Exa-MA contribution}\\
\midrule
\endfirsthead
\toprule
\textbf{Date} & \textbf{Announcement or event} & \textbf{Exa-MA contribution}\\
\midrule
\endhead
\midrule
\multicolumn{3}{r@{}}{\emph{Continued on next page}}\\
\endfoot
\bottomrule
\endlastfoot
12 January &
\href{https://www.hidalgo2.eu/webinar-on-global-challenges-and-built-environments-discover-4-hpc-applications-use-cases/}{HiDALGO2 and NCC France webinar} &
WP1 and WP7 presented Ktirio Urban Building for city-scale energy simulation.\\
14 January &
\href{https://exa-ma.numpex.fr/news/archive.html}{Website update} &
Expanded work package pages and software catalogue, with links to
reports and publications.\\
19--21 January &
\href{https://exa-ma.numpex.fr/news/2026/20260119-21-ag.html}{General Assembly, Aix-en-Provence} &
WP reviews, software sessions and a discussion with the Scientific
Advisory Board. Daria Hrebenshchykova became Exa-MA's YoungPEx representative.\\
2--5 February &
\href{https://exa-ma.numpex.fr/news/2026/20260202-05-transnumeriques.html}{Trans-Numériques, Rennes} &
Exa-MA contributed sessions on AI and HPC and on the software
ecosystem.\\
18--19 February &
\href{https://exa-ma.numpex.fr/news/archive.html}{AMD and NumPEx technical meeting, Saclay} &
WP1, WP3 and WP5 proposed topics for joint pilot studies with AMD.\\
7 May &
\href{https://exa-ma.numpex.fr/news/archive.html}{NumPEx Operational Committee, Paris} &
Exa-MA raised the deployment of common benchmarks on heterogeneous
systems.\\
26--28 May &
\href{https://aissai-hace.sciencesconf.org/?lang=en}{AISSAI--NumPEx HACE workshop, Toulouse} &
Emmanuel Franck (WP2) chaired the session on AI-based surrogates for PDEs.\\
8--12 June &
\href{https://majsc2026.pages.math.unistra.fr/program.html}{MAJSC workshop, Strasbourg} &
Lectures and practical sessions on JAX and automatic differentiation
for scientific computing.\\
11 June &
\href{https://exa-ma.numpex.fr/news/archive.html}{NumPEx annual review with the ANR} &
Exa-MA presented scientific results, collaboration and progress on the
project milestones.\\
22--26 June &
\href{https://exa-ma.numpex.fr/news/2026/20260622-hpddm-petsc-framework-webinar.html}{HPDDM/PETSc} (22),
\href{https://exa-ma.numpex.fr/news/2026/20260623-trust-framework-webinar.html}{TRUST} (23),
\href{https://exa-ma.numpex.fr/news/2026/20260625-freefem-framework-webinar.html}{FreeFEM++} (25) and
\href{https://exa-ma.numpex.fr/news/2026/20260626-samurai-framework-webinar.html}{Samurai} (26) &
Four software sessions on GPU preconditioning, Kokkos porting,
finite-element solver interfaces and adaptive mesh refinement.\\
7--9 September &
\href{https://aissai-jaxhack.sciencesconf.org/}{AISSAI--NumPEx JAX hackathon, IDRIS} &
Emmanuel Franck (WP2) taught an introductory JAX session for scientific computing.\\
6 November &
\href{https://numpex.org/4252/}{NumPEx position paper on agentic AI for science} &
Exa-MA contributed proposals on numerical validation, reproducibility
and benchmarking.\\
17 November &
\href{https://exa-ma.numpex.fr/news/2026/20261117-fem.html}{Forum Entreprises \& Mathématiques, Paris} &
Christophe Prud'homme moderated the HPC--AI panel on skills and careers.
Emmanuel Franck was a panellist.\\
\end{longtable}
\endgroup
